\chapter*{Preface}
\emph{Order and Chance} is a journey through three territories of mathematics, number theory, algebra and
probability, from the first principles to the questions that researchers ask today. Its destination is the place
where the three meet: the statistical behaviour of prime numbers, where the most rigid objects of mathematics turn
out to obey the laws of chance. Along the way every statement is proved. Nothing is asked to be taken on trust
that the series does not prove at some point, and the few foundations assumed here are named as such.

This first volume lays the ground, in eight chapters: sums, binomial coefficients and induction; complex numbers
and roots of unity; polynomials; inequalities; sequences, series and limits; integrals and Taylor's formulas;
linear algebra; and counting, with the first steps in probability. These are the tools every later tome uses, and
they are beautiful in their own right: the roots of unity turn a sum of binomial coefficients into a geometric
picture, a single inequality between means drives a dozen optimisation problems, and a determinant measures how a
linear map stretches volume.

Each chapter states the essentials briefly and then spends most of its pages on exercises, because mathematics is
learned by doing it. Exercises are graded \textbf{$\star$} for a warm-up, \textbf{$\star\star$} for a good afternoon
and \textbf{$\star\star\star$} for gems, and every one has a hint and a complete solution at the end of its chapter. Each solution
begins with \emph{How to find it}, which explains what to notice and why the method is natural, before the
\emph{Write-up}: the aim is to learn how a proof is found, not only how it reads. ``Pause and think'' boxes ask
you to guess before you read on; the guess is the point, not whether it is right.

Every theorem, proposition and lemma is proved. The only facts taken without proof are the foundations: the
axioms of $\N$ and of $\R$, the existence of the integral of a continuous function, the intermediate value
theorem, the extreme values of continuous functions on closed bounded sets, Heine's theorem on uniform
continuity and the mean value theorem that follows from them, and the basic properties of $\exp$, $\ln$, $\cos$,
$\sin$ and $\pi$. Tome~6 constructs the real numbers and proves every fact in this list after the axioms. Results quoted only for their
interest appear as remarks, with a reference to where they are proved.

Each chapter has a companion Python notebook, in which the identities of the chapter are tested numerically and
put to work, often on a problem from computer science.
